\documentclass[11pt]{article}
\usepackage[margin=1in]{geometry}
\usepackage[hidelinks]{hyperref}
\usepackage{xurl}
\usepackage{enumitem}

\title{Congestion Sub-Series --- Index (Unified Source-Only Archive)}
\author{}
\date{Unified archive v121 (2026-02-24)}

\begin{document}
\maketitle

\paragraph{Scope.}
This note indexes the Congestion sub-series and records where historical ``split supplement'' content was absorbed into the unified source-only papers.

\paragraph{Boundary.}
This file owns only the family map: a reading order, one-sentence role split across the congestion papers, and a continuity pointer into the archive-wide series index. It does not own the congestion witness/accountant contract, the PSC-Q mechanism packet, the verifier bundle, evaluator notarization, runtime conformance receipts, or public release lineage.

\paragraph{Imports.}
None; this file is purely navigational.

\paragraph{Exports.}
A three-paper primary reading order, one optional addendum pointer, a short absorbed-supplement continuity note, and a continuity pointer into the series-wide index (\path{../../../index/SERIES_INDEX.tex}).

\paragraph{Earliest sufficient stop.}
If a reader can identify the witness/accountant paper, the mechanism packet paper, the verifier-bundle paper, and the optional supporting-evidence addendum, then this note has done enough.

\paragraph{Artifact policy.}
Source-only; plots and regression vectors referenced by the papers are available on request.



\section*{What this sub-series contains}
The congestion sub-series has a three-paper primary pipeline plus one optional supporting-evidence addendum. See \path{../../../index/SERIES_INDEX.tex} for the continuity map.
\begin{itemize}[leftmargin=1.2em]
\item \textbf{Congestion-EQ} (\path{paper1\_congestion\_eq/paper.tex}): a TV-tier contract and accounting rules for congestion witnesses (RTT, queueing delay, loss/ECN marks).
\item \textbf{PSC-Q} (\path{paper2\_psc\_q/paper.tex}): the declared mechanism packet (public service calendars, slack-harvesting substitution padding, dither/expiry knobs, utilization range, and observer assumptions) for meeting Congestion-EQ budgets in compiler-friendly systems.
\item \textbf{W-Congestion-EQ} (\path{paper3\_w\_congestion\_eq/paper.tex}): the congestion-specific verifier bundle (canonical fields, replay checks, checker outputs, and exported digests) for published congestion-budget claims.
\item \textbf{PSC-Q (Addendum)} (\path{paper2A\_psc\_q\_microstudies\_addendum/paper.tex}; optional companion): a supplementary evidence packet with omitted-figure placeholders, short microstudy summaries, and implementation cautions for PSC-Q operators.
\end{itemize}

\section*{Notes on tightness and ``supplement'' material}
The historical split archive shipped separate ``supplement'' PDFs with additional plots and regression vectors.
In this unified archive, we keep a \emph{single-file} rule: each paper is one \texttt{.tex} file.
Operator-facing microstudies and regression-vector descriptions are preserved either as short sections inside the main papers or, for PSC-Q, as the narrow supporting-evidence addendum. Large plots, scripts, and JSON evidence remain available on request as an external artifact bundle.

\section*{Publication posture}
This note is intentionally a maintainer-facing crosswalk rather than a theorem, certificate, or release paper. It is useful for waking up inside the archive and for preserving the congestion-family split after the old supplement material was reabsorbed, but it is not yet a strong public entrypoint on its own. The conservative default for this pass is therefore to keep it out of Candidate and record an explicit Hold instead of treating it like a releaseable research paper.

\end{document}
